\section{Field definition}
\label{sec:fielddef}

The computed quantity is a single linear integral. For a bunch of source charges
at positions $\mathbf{x}_p$ with signed charge $q_p$, and a shared
dimensionless velocity $\mathbf{u}=\mathbf{p}/(mc)$, the field at a staggered
observation point $\mathbf{x}_{\mathrm{obs}}$ is the TSC-convolved
uniformly-moving point-charge field
\begin{equation}
\label{eq:Fc}
F_c(\mathbf{x}_{\mathrm{obs}})
=
\sum_p q_p \int_{[-1.5,1.5]^3}
W(\mathbf{s})\,
K_{\mathrm{point},c}\!\bigl(\mathbf{x}_{\mathrm{obs}}
- \mathbf{x}_p - \mathbf{s}\,\mathbf{cell},\ \mathbf{u}\bigr)
\,d\mathbf{s}^3 ,
\end{equation}
where $W(\mathbf{s})=W_x(s_x)W_y(s_y)W_z(s_z)$ is the separable TSC shape
(Section~\ref{sec:method}) and $K_{\mathrm{point}}$ is the per-unit-charge
field of one uniformly moving charge. Because the integral is linear in the
source charge, the sum over particles is exact; the whole design is a fast
evaluation of this one integral.

\subsection{The moving point kernel}

The kernel is the present-time Li\'enard--Wiechert field of a uniformly moving
charge \citep{S1-spec}. For constant velocity the field points to the
instantaneous position. In the lab frame, with
$\mathbf{R}=\mathbf{x}_{\mathrm{obs}}-\mathbf{x}_{\mathrm{src}}$,
\begin{align}
\gamma &= \sqrt{1+|\mathbf{u}|^2}, &
\boldsymbol{\beta} &= \mathbf{u}/\gamma, &
\mathrm{invg}^2 &= 1/(1+|\mathbf{u}|^2) = 1-\beta^2,\\
D &= |\mathbf{R}|^2\,\mathrm{invg}^2 + (\mathbf{R}\!\cdot\!\boldsymbol{\beta})^2,\\
\mathbf{E}_n &= \mathrm{invg}^2\,\mathbf{R}/D^{3/2}, &
\mathbf{B}_n &= \boldsymbol{\beta}\times \mathbf{E}_n .
\end{align}
For $\mathbf{u}=\mathbf{0}$ exactly, $\mathbf{E}_n=\mathbf{R}/|\mathbf{R}|^3$
and $\mathbf{B}_n=\mathbf{0}$; no $\varepsilon$ is added to any denominator.
The kernel is homogeneous of degree $-2$ in $\mathbf{R}$. A bunch with a
velocity spread needs one convolution per velocity class
(Section~\ref{sec:multivel}); the validated scope is a single shared
$\mathbf{u}$.

\subsection{The unit chain}

The reference normalizes the per-unit-charge field as
$E_n = 4\pi\varepsilon_0\,E/q$ and $B_n = 4\pi\varepsilon_0 c\,B/q$, so
$B_n=\boldsymbol{\beta}\times E_n$ \citep{S1-spec}. The mesh produces the
per-unit-charge field with $\mathbf{R}$ in metres, already summed over the
deposited charge $q_p$ in coulombs. The chain to physical SI is
\begin{equation}
E_{\mathrm{SI}}[c] = F_c[c]\,h^3/(4\pi\varepsilon_0),\quad
B_{\mathrm{SI}}[c] = F_c[c]\,h^3/(4\pi\varepsilon_0 c),
\end{equation}
with $\varepsilon_0=8.8541878128\times10^{-12}$ F/m and
$c=299792458$ m/s. The $h^3$ factor (Section~\ref{sec:method}) is the Riemann
volume of the sampled continuum integral \citep{S1-spec}.

\subsection{The TSC shape}

$W$ is the quadratic B-spline, with argument in cell units, support
$[-1.5,1.5]$ per axis, and unit integral per axis
\citep{S1-spec}:
\begin{equation}
W(x) =
\begin{cases}
3/4 - x^2 & |x|<1/2,\\
\frac{1}{2}(3/2-|x|)^2 & 1/2\le |x|<3/2,\\
0 & \text{otherwise.}
\end{cases}
\end{equation}
